\begin{definition}[Types and observable equality]\label{def:registers}
A register list $L$ is a finite ordered list of distinct wire names, with
$\cH_L=(\C^2)^{\otimes |L|}$. Subregisters are ordered subsets; juxtaposition
of disjoint lists denotes tensor product. Products of operators act from
right to left. A width-$w$ value is a binary string, also interpreted as
an unsigned integer in $\{0,\ldots,2^w-1\}$; $\oplus$ is bitwise XOR.
A density operator is a positive operator of trace one. A channel is a
completely positive trace-preserving linear map. An instrument is a finite
family of completely positive trace-nonincreasing maps summing to a channel.
For channels with the same input and output, define
\[
 d_{\rm out}(\Phi,\Psi)=
 \sup_{R,\rho}\|((\Phi-\Psi)\otimes\id_R)(\rho)\|_1,
\]
where $R$ ranges over finite reference registers and $\rho$ over density
operators on the input and $R$; $\|A\|_1=\tr\sqrt{A^\dagger A}$.
An instrument is compared as a channel retaining its orthogonal outcome
labels. The operator norm is $\|A\|=\sup_{\|v\|=1}\|Av\|$.
All bounds and circuit descriptions are independent of the unknown input.
\end{definition}

\begin{definition}[Machine 3 and fresh names]\label{def:circuit}
Machine 3 is the finite quantum circuit model with labelled gates
\[
 H=2^{-1/2}\begin{pmatrix}1&1\\1&-1\end{pmatrix},\quad
 T=\operatorname{diag}(1,e^{i\pi/4}),\quad
 \operatorname{CNOT}\ket{a,b}=\ket{a,b\oplus a}.
\]
Write $\mathcal G=\{H,T,\operatorname{CNOT}\}$,
$\mathcal G^{\pm}=\mathcal G\cup\{T^\dagger\}$,
$Z=T^4$, $X=HZH$, and $S=T^2$.
Exactly two additional register transformations are allowed:
\[
 A_B:\cH_L\longrightarrow\cH_{LB},\quad
 A_B\ket\psi=\ket\psi\ket{0^{|B|}},\qquad
 D_B(\rho)=\tr_B\rho\quad(B\subseteq L).
\]
The names in $B$ in $A_B$ have never occurred earlier in the description.
Drop removes names from the accessible list permanently; later allocations
use new names. The circuit data designate an accessible output list $O$.
In the retained interpretation, a dropped factor is moved to an inaccessible
list $E$ and kept, with no subsequent operation on it. Thus the denotation
is first an isometry $V:\cH_{L_{\rm in}}\to\cH_{OE}$ and then the channel
$\rho\mapsto\tr_E(V\rho V^\dagger)$.
An optional ideal alphabet adds the actual matrices
$R_z(\theta)=\operatorname{diag}(e^{-i\theta/2},e^{i\theta/2})$.
An angle is supplied by a finite description together with a procedure
that, on input $b$, outputs dyadic approximations to the four complex
entries of $R_z(\theta)$ within $2^{-b}$ (the entry procedure of
Theorem~\ref{eqv:rotations}, whose cost is charged there). A description
without such a procedure does not define an ideal gate in this paper.
Their controlled versions use these matrices including their phases;
$T$ is not identified with $R_z(\pi/4)$ under a control.

Terminal computational-basis observation is part of the circuit model,
not an additional register transformation. Before choosing $O$, let $A$
be the final accessible list. The data may designate an ordered list
$(i_1,\ldots,i_b)$ of wires in $A$ to measure (repetitions allowed).
For $y\in\{0,1\}^b$ set
$P_y=\prod_{j=1}^b(\proj{y_j}_{i_j}\otimes\id)$.
Let $E_0$ be the list of previously dropped wires, disjoint from $A$.
If $V_0:\cH_{L_{\rm in}}\to\cH_{AE_0}$ is the retained circuit
isometry before observation, the
instrument on $O\subseteq A$ is
\[
 \rho\longmapsto\tr_{(AE_0)\setminus O}
       (P_yV_0\rho V_0^\dagger P_y),\qquad y\in\{0,1\}^b,
\]
where projectors act as identity on $E_0$. Its retained form first
copies each $i_j$ into
a fresh zero wire $Y_j$ by CNOT, then copies $Y$ into a fresh zero
$E_Y$. Selecting $OY$ and tracing its complement gives the displayed
instrument with orthogonal classical labels $y$. Without observation
the earlier output isometry notation is unchanged; with observation
it has output $OY$ and the complementary retained environment.
\end{definition}

